\documentclass[a4paper,12pt]{report}

% Paquetes necesarios
\usepackage[utf8]{inputenc}   % Codificación UTF-8
\usepackage{fontspec} % Paquete para manejar fuentes específicas
\usepackage[spanish]{babel}
\usepackage[style=apa, backend=biber]{biblatex}
\usepackage{geometry}         % Configuración de márgenes
\usepackage{graphicx}
\usepackage{setspace}         % Configuración de interlineado
\usepackage{fancyhdr}         % Configuración de encabezados y pie de página
\usepackage{titlesec} % Para personalizar títulos de capítulos
\usepackage{hyperref} % Opcional: Para hipervínculos en el índice



\setmainfont{Arial}   % Configura Arial como la fuente principal
\geometry{
    top=2.5cm,
    bottom=2.5cm,
    left=4cm,
    right=2.5cm
}
\pagestyle{fancy}
% Configuración del interlineado y espacio entre párrafos
% \setlength{\parskip}{2em}  % Espacio doble entre párrafos
% \setlength{\parindent}{0em} % Sin sangría en párrafos
\onehalfspacing            % Interlineado de 1.5
% Configuración del pie de página
\fancyhf{}
\fancyfoot[R]{\thepage}    % Número de página en la esquina inferior derecha
\fancypagestyle{plain}{
    \fancyhf{}
    \fancyfoot[R]{\thepage}
}

% Configuración para personalizar capítulos sin la palabra "Capítulo"
\newcommand{\customchapter}[1]{
    \chapter*{#1} % Capítulo sin "Capítulo" y sin numeración
    % \addcontentsline{toc}{chapter}{#1} % Agrega al índice sin numerar
}

% Configuración para personalizar títulos de capítulos específicos
\newcommand{\nochapterlabel}{
  \titleformat{\chapter}[hang]
    {\normalfont\huge\bfseries}
    {\thechapter.}{1em}{} % Quita "Capítulo" pero deja el número
}


\addbibresource{bibliography.bib} % Archivo .bib con las referencias



\begin{document}

% Incluir la portada
% Archivo: portada.tex
\begin{titlepage}
    \begin{center}
        {\fontsize{15pt}{10pt}\selectfont\textbf{UNIVERSIDAD AUTÓNOMA "GABRIEL RENE MORENO"}} \\
        {\fontsize{15pt}{10pt}\selectfont\textbf{FACULTAD DE INGENIERÍA EN CIENCIAS DE LA COMPUTACIÓN Y TELECOMUNICACIONES}} \\
        \vspace{1cm}
        % Logo
        \includegraphics[width=0.8\textwidth]{images/logo_soe.png} \\
        
        % Título de la institución
        {\fontsize{15pt}{10pt}\selectfont\textbf{MAESTRÍA EN CIENCIA DE DATOS E INTELIGENCIA ARTIFICIAL}} \\

        % Título del trabajo
        \vspace{1cm}
        {\fontsize{15pt}{5pt}\selectfont\textbf{Implementar un Chatbot Semántico Multiperfil para Acceso Inteligente a Documentación Técnica y de Negocio en Plataformas Confluence del Sector Aeronáutico en Airnguru S.A. para la Gestión 2025}} \\
        \vspace{1cm}
        
        \vspace{2cm}
        \begin{flushright}
        {\fontsize{14pt}{22pt}\selectfont\textbf{Autor}} \\
        {\fontsize{14pt}{22pt}\selectfont{Ing. Darlyn Bravo Peña}} \\
        \end{flushright}
        
        
        
        % Ciudad y fecha
        \vfill
        \textbf{Santa Cruz, Bolivia} \\
    \end{center}
    \end{titlepage}
     % Reemplaza con % Archivo: portada.tex
\begin{titlepage}
    \begin{center}
        {\fontsize{15pt}{10pt}\selectfont\textbf{UNIVERSIDAD AUTÓNOMA "GABRIEL RENE MORENO"}} \\
        {\fontsize{15pt}{10pt}\selectfont\textbf{FACULTAD DE INGENIERÍA EN CIENCIAS DE LA COMPUTACIÓN Y TELECOMUNICACIONES}} \\
        \vspace{1cm}
        % Logo
        \includegraphics[width=0.8\textwidth]{images/logo_soe.png} \\
        
        % Título de la institución
        {\fontsize{15pt}{10pt}\selectfont\textbf{MAESTRÍA EN CIENCIA DE DATOS E INTELIGENCIA ARTIFICIAL}} \\

        % Título del trabajo
        \vspace{1cm}
        {\fontsize{15pt}{5pt}\selectfont\textbf{Implementar un Chatbot Semántico Multiperfil para Acceso Inteligente a Documentación Técnica y de Negocio en Plataformas Confluence del Sector Aeronáutico en Airnguru S.A. para la Gestión 2025}} \\
        \vspace{1cm}
        
        \vspace{2cm}
        \begin{flushright}
        {\fontsize{14pt}{22pt}\selectfont\textbf{Autor}} \\
        {\fontsize{14pt}{22pt}\selectfont{Ing. Darlyn Bravo Peña}} \\
        \end{flushright}
        
        
        
        % Ciudad y fecha
        \vfill
        \textbf{Santa Cruz, Bolivia} \\
    \end{center}
    \end{titlepage}
     si prefieres que se trate como una sección independiente

% Table of contents
\newpage
\tableofcontents
\newpage
\listoffigures
\newpage
\listoftables
\newpage
\begin{center}
    {\LARGE \textbf{INTRODUCCIÓN}} % Título centrado y en negrita
\end{center}

La estimación de costos en el desarrollo de software es un desafío inherente a la naturaleza intangible de este producto. En un mundo donde la economía digital impulsa la innovación y la competitividad, la incapacidad de calcular con precisión los costos de desarrollo puede llevar a pérdidas financieras, proyectos fallidos y relaciones comerciales insatisfactorias. A nivel global, empresas de todos los tamaños enfrentan la necesidad de establecer modelos más eficientes y estandarizados para valorar su trabajo, especialmente en un mercado donde la externalización y la globalización son cada vez más comunes. Este tema adquiere mayor relevancia 
con el auge de las metodologías ágiles y la inteligencia artificial, que prometen pero también complican la predicción de costos.

En Bolivia, donde el sector tecnológico está en crecimiento, el problema de la estimación de costos en el desarrollo de software adquiere un matiz especial. La falta de modelos estandarizados y herramientas locales para valorar los servicios de desarrollo de software lleva a una alta dependencia de las negociaciones subjetivas, lo que puede perjudicar tanto a los desarrolladores como a los clientes. Esto es especialmente importante en un contexto donde muchas empresas están comenzando a digitalizarse y dependen del software para competir en un mercado nacional e internacional. Abordar este desafío no solo contribuirá al fortalecimiento del sector tecnológico en Bolivia, sino que también permitirá que los desarrolladores y empresas logren acuerdos más justos y sostenibles, impulsando así el crecimiento de la economía digital en el país.

El tema está ubicado en el area 
Ciencias Generales de la Computación Aplicadas y la 
línea Ingeniería de software con un eje de 
Aspectos económicos y de negocio del proceso de desarrollo del software porque la estimación de costos afecta la 
rentabiliad de los proyectos del software, tanto para los desarrolladores
como para los clientes.
\section{Antecedentes}
La estimación precisa de costos en el desarrollo de software es un desafío constante en la ingeniería de software, especialmente en entornos ágiles. Técnicas como los \textit{Puntos de Historia} y el \textit{Planning Poker} han sido desarrolladas para abordar este problema, permitiendo a los equipos estimar el esfuerzo necesario para completar historias de usuario de manera colaborativa 
y basada en consenso \parencite{asana_planning_poker}
Los \textit{Puntos de Historia} son una unidad de medida abstracta que refleja el esfuerzo relativo, la complejidad y el riesgo asociados a una historia de usuario. Esta técnica permite a los equipos evaluar tareas en función de su tamaño relativo, en lugar de estimar en unidades de tiempo, lo que facilita 
la planificación y priorización en entornos ágiles \parencite{asana_story_points}.

Por otro lado, el \textit{Planning Poker} es una técnica de estimación que combina la opinión de todos los miembros del equipo para alcanzar una estimación consensuada. Durante una sesión de \textit{Planning Poker}, cada miembro selecciona de manera independiente una carta con el valor que considera adecuado para una tarea específica. Posteriormente, las discrepancias se discuten hasta alcanzar un consenso, promoviendo 
la colaboración y la precisión en las estimaciones \parencite{asana_planning_poker}.

A pesar de su efectividad en proyectos ágiles, estas técnicas enfrentan desafíos, como la subjetividad inherente al proceso y la falta de estandarización en algunos equipos. Sin embargo, su uso ha demostrado ser fundamental para mejorar la precisión de las estimaciones y la planificación 
en proyectos de software de diversa escala \parencite{asana_planning_poker}.

\section{Presentación de las Problemáticas}

El enfoque ágil, específicamente mediante técnicas como los \textit{puntos de historia} y el \textit{Planning Poker}, busca mejorar la precisión en las estimaciones promoviendo la colaboración entre equipos y utilizando métricas iterativas y relativas. Sin embargo, estas metodologías enfrentan problemáticas clave, como:

\begin{itemize}
    \item \textbf{Subjetividad en las estimaciones:} Las valoraciones basadas en puntos de historia pueden variar significativamente según la experiencia y perspectiva de los miembros del equipo.
    \item \textbf{Dependencia de la experiencia del equipo:} La efectividad de \textit{Planning Poker} depende de un equipo con experiencia previa en proyectos similares, lo que limita su aplicabilidad en equipos menos maduros.
    \item \textbf{Dificultad en la estandarización:} Los puntos de historia y \textit{Planning Poker} no siempre son fácilmente comparables entre equipos o proyectos, lo que complica la evaluación del desempeño y la predicción futura.
    \item \textbf{Falta de datos históricos:} En proyectos nuevos o innovadores, la ausencia de datos históricos dificulta establecer referencias confiables para las estimaciones.
\end{itemize}


\section{Situación Problemática}

La estimación de costos y tiempos en el desarrollo de software es un componente crítico para garantizar el éxito de un proyecto. En entornos de trabajo tradicionales, estas estimaciones suelen realizarse utilizando métodos rígidos y lineales que no se adaptan fácilmente a los cambios frecuentes en los requerimientos. Sin embargo, con la creciente adopción de metodologías ágiles, han surgido técnicas como los \textit{puntos de historia} y el \textit{Planning Poker}, que buscan proporcionar mayor flexibilidad y precisión en estas estimaciones.

A pesar de su potencial, estas técnicas presentan múltiples desafíos en su implementación. Por ejemplo, las estimaciones realizadas con puntos de historia suelen ser subjetivas, dependiendo de la experiencia y el criterio de los miembros del equipo. Además, el uso de \textit{Planning Poker} requiere equipos bien entrenados y con experiencia, lo que puede dificultar su aplicación en proyectos con equipos nuevos o poco maduros.

Otro problema significativo es la falta de estandarización en estas técnicas. Los puntos de historia no son fácilmente comparables entre diferentes proyectos o equipos, lo que complica la evaluación del desempeño y la generación de predicciones confiables. Asimismo, en proyectos innovadores o en fases iniciales, la ausencia de datos históricos dificulta establecer referencias para realizar estimaciones precisas.

Estas problemáticas subrayan la importancia de analizar y evaluar la efectividad de los \textit{puntos de historia} y el \textit{Planning Poker} como enfoques para la estimación de costos y tiempos en entornos ágiles. Comprender sus limitaciones y áreas de mejora puede contribuir al desarrollo de estrategias más efectivas para la planificación y ejecución de proyectos de software.

\section{Pregunta de Investigación}
¿Cómo impacta el uso de puntos de historia y Planning Poker en la precisión y eficiencia de la estimación de costos en proyectos de desarrollo de software bajo metodologías ágiles?
\section{Justificación Teórica}

La estimación de costos y tiempos en el desarrollo de software es una actividad esencial dentro de la ingeniería de software, dado su impacto en la planificación y ejecución de proyectos. Desde una perspectiva teórica, las metodologías ágiles han emergido como una respuesta a los retos planteados por los enfoques tradicionales, promoviendo principios como la adaptabilidad, la iteración y la colaboración constante entre los miembros del equipo.

Dentro de este marco, los \textit{puntos de historia} y el \textit{Planning Poker} representan herramientas clave basadas en teorías de estimación relativa y trabajo en equipo. Estas técnicas permiten abordar problemas complejos mediante la descomposición de tareas y la construcción de consensos. Sin embargo, su uso plantea interrogantes teóricas importantes, como la subjetividad inherente a las estimaciones, la dependencia de la experiencia del equipo y la falta de modelos estandarizados que permitan extrapolar resultados entre proyectos.

Estudiar estas técnicas desde una perspectiva teórica permite evaluar su validez y utilidad en comparación con enfoques tradicionales de estimación. Además, ayuda a construir un marco conceptual que respalde el diseño e implementación de metodologías ágiles en diversos contextos, enriqueciendo el campo de la gestión de proyectos de software.

\section{Justificación Práctica}

En el ámbito práctico, las metodologías ágiles han ganado una amplia aceptación en la industria del software debido a su capacidad para adaptarse a cambios rápidos en los requerimientos y mejorar la satisfacción del cliente. Sin embargo, la estimación de costos y tiempos sigue siendo un desafío significativo, ya que los errores en esta etapa pueden conducir a sobrecostos, retrasos y la insatisfacción de los clientes.

El uso de \textit{puntos de historia} y \textit{Planning Poker} en la estimación ágil ofrece ventajas prácticas, como fomentar la participación del equipo, promover el consenso y facilitar la planificación iterativa. No obstante, estas técnicas presentan problemas prácticos, como la dificultad de estandarización entre equipos, la variabilidad en las estimaciones debido a diferencias de experiencia y la falta de herramientas para medir su efectividad en términos reales.

Este estudio busca abordar estos problemas, proporcionando un análisis crítico de estas técnicas en entornos reales de desarrollo de software. Sus resultados podrán ofrecer a los equipos de desarrollo recomendaciones prácticas para mejorar la precisión de sus estimaciones, optimizar recursos y aumentar la efectividad en la entrega de proyectos ágiles.

\section{Objetivo General}

Analizar la efectividad de las técnicas de estimación ágiles, específicamente los \textit{puntos de historia} y el \textit{Planning Poker}, en la estimación de costos en proyectos de desarrollo de software, evaluando su impacto en la precisión, eficiencia y planificación de recursos dentro de equipos ágiles.

\section{Objetivos Específicos}

\begin{itemize}
    \item \textbf{Identificar} las principales características y fundamentos teóricos de las técnicas de estimación ágiles, como los \textit{puntos de historia} y el \textit{Planning Poker}.
    \item \textbf{Evaluar} la precisión de las estimaciones realizadas con \textit{puntos de historia} y \textit{Planning Poker}, comparándolas con métodos tradicionales de estimación.
    \item \textbf{Determinar} cómo la experiencia y composición del equipo afectan la efectividad de estas técnicas de estimación en proyectos ágiles.
    \item \textbf{Analizar} las principales limitaciones y desafíos asociados con el uso de estas herramientas en entornos de desarrollo de software.
    \item \textbf{Proponer} recomendaciones basadas en los hallazgos del análisis para optimizar la implementación de \textit{puntos de historia} y \textit{Planning Poker} en equipos ágiles.
\end{itemize}

\chapter{Delimitación o Alcance}

\section{Delimitación o Alcance Espacial}

El estudio se llevará a cabo en organizaciones y empresas dedicadas al desarrollo de software que aplican metodologías ágiles. La investigación se centrará en proyectos realizados en la ciudad de Santa Cruz de la Sierra, Bolivia, dado que esta región ha experimentado un crecimiento considerable en la industria tecnológica y en la adopción de metodologías ágiles para la planificación y ejecución de proyectos de software. 

El análisis se enfocará en equipos de desarrollo que utilizan herramientas específicas como \textit{puntos de historia} y \textit{Planning Poker}, permitiendo una evaluación contextualizada de su aplicación en el ámbito local. Aunque el alcance geográfico se limita a Santa Cruz, los resultados podrán servir como referencia para organizaciones con características similares en otras regiones.

\section{Alcance Temporal}

El alcance temporal del estudio abarca el análisis de proyectos de desarrollo de software realizados entre los años 2020 y 2024. Este período se selecciona debido a la creciente adopción de metodologías ágiles en la industria del software durante estos años, impulsada por la necesidad de adaptarse a entornos dinámicos y cambiantes. Además, este rango de tiempo permite evaluar cómo las técnicas de \textit{puntos de historia} y \textit{Planning Poker} han evolucionado y se han implementado en proyectos recientes, asegurando la relevancia y actualidad de los datos recopilados.


\newpage
\printbibliography

\end{document}
